\lecture{14}{Wed 06 Nov 2024 10:47}{Insanity and More Infinite Dimensional Systems}
\textbf{ASIDE}: Let $\mathbf{F}=(x,y,z)$ be the solution to some 3 dimensional system of differential equations that orbits $\mathbb{R}^3$. Define a new $\xi \eta $-plane that is orthogonal to $\mathbf{F}$ at some point. As $t$ increases, the equation will loop back around and continue to intersect the $\xi \eta $-plane, meaning we can define a discrete dynamical system $P(\xi ,\eta )$ as simply the points $\mathbf{F}$ cuts the $\xi \eta $-plane.

Change in notation: $v(x)$ is now called $\phi (x)$. Recall the question: are there solutions $\phi (x)$ such that $\phi (0)=\phi (\ell )=0$, where $\phi $ solves
\[
	\lambda \phi (x)=\frac{\mathrm{d}^2\phi }{\mathrm{d}x^2} 
?\]
We claim all nonzero solutions have $\lambda \in\mathbb{R}$. Suppose $\lambda ,\phi (x)$ solve the problem. Note that if we take the complex conjugate, we find  $\overline{\lambda} $ and $\overline{\phi } $ also solve the equation. If $\phi (0)=\phi (L)=0$, then $\overline{\phi (0)} =\overline{\phi (L)} =0$. Multiply through by $\overline{\phi  } $ :
\[
	\overline{\phi }\cdot  \frac{\mathrm{d}^2\phi }{\mathrm{d}x^2} =\lambda \left| \phi (x) \right| ^2
.\]
Integrating both sides, it follows that
\[
	\int_{0}^{L} \overline{\phi (x)} \cdot \frac{\mathrm{d}^2\phi }{\mathrm{d}x^2}  \,\mathrm{d} x=\lambda \int_{0}^{L} \left| \phi (x) \right| ^2 \,\mathrm{d} x
.\]
Using IBP on the left derivative, we have
\[
	\left.\overline{\phi (x)} \frac{\mathrm{d}\phi }{\mathrm{d}x}\right|_{0}^{L} -\int_{0}^{L} \frac{\mathrm{d} \overline{\phi  }}{\mathrm{d}x} \frac{\mathrm{d}\phi }{\mathrm{d}x}  \,\mathrm{d} x = \left. \overline{\phi (x)} \frac{\mathrm{d}\phi }{\mathrm{d}x} \right|_{0}^{\ell }-\int_{0}^{L} \left| \frac{\mathrm{d}\phi }{\mathrm{d}x}  \right| ^2 \,\mathrm{d} x
.\]
Using our boundary terms, we know
\[
	\left.\overline{\phi (x)} \frac{\mathrm{d}\phi }{\mathrm{d}x}\right|_{0}^{L} =0
.\]
Our integrand is thus
\[
	\int_{0}^{L} \left| \frac{\mathrm{d}\phi }{\mathrm{d}x}  \right|  \,\mathrm{d} x
.\]
Recalling our previous integral and solving for $\lambda $,
\[
	\lambda =\frac{-\int_{0}^{L} \left| \frac{\mathrm{d}\phi }{\mathrm{d}x}  \right| ^2 \,\mathrm{d} x}{\int_{0}^{L} \left| \phi (x) \right| ^2 \,\mathrm{d} x}\in\mathbb{R}
.\]

We now claim that if $(\lambda _1,\phi _1),(\lambda _2,\phi _2)$ are solutions of our problem and $\lambda _1\neq \lambda _2$, then $\left<\phi_1,\phi _2 \right> =0$, where we recall that for functions $\phi ,\psi $, the inner product is defined as
\[
	\left<\phi ,\psi  \right> =\int_{0}^{L} \phi (x)\psi (x) \,\mathrm{d} x
.\]
Consider the inner product
\[
	\lambda _1\left<\phi_1,\phi _2\right> =\int_{0}^{L} \lambda_1\phi_1\phi _2 \,\mathrm{d} x
.\]
Recall that we found
\[
	\frac{\mathrm{d}^2u}{\mathrm{d}x^2} =\lambda \phi 
.\]
Thus,
\[
	\int_{0}^{L} \frac{\mathrm{d}^2\phi_1 }{\mathrm{d}x^2} \phi _2 \,\mathrm{d} x
.\]
Using IBP,
\[
	\left. \frac{\mathrm{d}\phi }{\mathrm{d}x} \phi _2(x)\right|_{0}^{L}-\int_{0}^{L} \frac{\mathrm{d}\phi _1}{\mathrm{d}x} \frac{\mathrm{d}\phi _2}{\mathrm{d}x}  \,\mathrm{d} x
.\]
Recalling that $\phi_1,\phi _2$ solve the equation, the boundary terms go to 0. Using IBP again, we have
\[
	\left. -\phi _1(x)\frac{\mathrm{d}\phi _2}{\mathrm{d}x} \right|_{0}^{L}+\int_{0}^{L} \phi _1(x)\frac{\mathrm{d}^2\phi _2}{\mathrm{d}x^2}  \,\mathrm{d} x
.\]
The boundary terms once again go to 0, so plugging in $\lambda _2\phi _2=\frac{\mathrm{d}^2\phi_2 }{\mathrm{d}x} $, we have
\[
	\int_{0}^{L} \phi _1(x)\lambda_2\phi _2(x) \,\mathrm{d} x=\lambda \left<\phi _1,\phi _2 \right>
.\]
Therefore,
\[
	\lambda _1 \left<\phi _1(x),\phi _2(x) \right> =\lambda _2\left<\phi _1(x),\phi _2(x) \right>
.\]
Recall that we suppose $\lambda _1\neq \lambda _2$. Therefore,
\[
	\left<\phi _1(x),\phi _2(x) \right> =0
.\]

Now, we will solve for the eigenfunctions of this problem. Since $\lambda $ is negative and real, write $\lambda =-\omega ^2$. Then
\[
	\frac{\mathrm{d}^2\phi }{\mathrm{d}x^2} =-\omega ^2 \phi (x)
.\]
We know we will have two solutions, one in $\sin $ and one in $\cos $. We could also solve this with an exponential using euler's formula, but that's boring. We have
\[
	\phi (x)=A \cos (\omega x)+V\sin (\omega x)
.\]
We still have to solve the boundary conditions, however. Letting $x=0$, we have $\phi (0)=A$. Therefore, $A=0$. Set $\phi (\ell )=0$. We have
\[
	\phi (\ell )=B\sin (\omega \ell )=0
.\]
We must have $B\neq 0$, meaning $\omega \ell =n\pi $ for some $n\in\mathbb{N}$. Thus, we have as a solution $\phi (x)=B\sin (\omega x)$, provided
\[
	\omega =\frac{n\pi }{\ell }
.\]
Therefore, we have infinite eigenvalues
\[
	\lambda_n = -\omega_n^2
,\]
and infinite associated eigenfunctions
\[
	\phi _n(x)=B_n\sin \left( \frac{n\pi x}{\ell } \right) 
.\]
Recall the original function
\[
	u(x,t)=e^{-\lambda t}\phi (x)
.\]
We have
\[
	u_n(x,t)=e^{-\left( \frac{n\pi }{\ell } \right) ^2t}\sin \left( \frac{n\pi x}{\ell } \right) 
.\]
Although these solutions satisfy the boundary conditions and the heat equation, they don't necessarily satisfy the initial conditions \textbf{unless} the initial heat in the wire happens to be some sin function. However, even in 2D systems the soltuions generally didn't solve the IVPs. So we took a linear combination of solutions and chose the appropriate constants. Similarly, we take a linear combination of all eigensolutions as our general solution:
\[
	u(x,t)=\sum_{n=1}^{\infty} B_n\exp \left( -\left(\frac{n\pi }{\ell }\right)^2t \right) \sin \left( \frac{n\pi x}{\ell } \right) 
.\]
In order to satisfy the initial conditions $u(x,0)$, we would need
\[
	\sum_{n=1}^{\infty} B_n \sin \left( \frac{n\pi x}{\ell } \right) \exp (0)=\sum_{n=1}^{\infty} B_n \sin \left( \frac{n\pi x}{\ell } \right) =u_0(x)
.\]
This is a Fourier series. Recall that conditions exist such that this is true. Multiply both sides by $\sin (m\pi x /\ell )$. We have
\[
	\sum_{n=1}^{\infty} \sin \left( \frac{m\pi x}{\ell } \right) B_n\sin \left( \frac{n\pi x}{\ell } \right) =u_0 \sin \left( \frac{m\pi x}{\ell } \right) 
.\]
We will integrate both sides from $0$ to $\ell $.
\[
	\int_{0}^{\ell }\sum_{n=1}^{\infty} \sin \left( \frac{m\pi x}{\ell } \right) B_n\sin \left( \frac{n\pi x}{\ell } \right)\mathrm{d} x =\int_{0}^{L}u_0 \sin \left( \frac{m\pi x}{\ell } \right) \mathrm{d}x 
.\]
We will interchange the sum with the integral because that works for analytic series.
\[
	\sum_{n=1}^{\infty}\left( B_n \int_{0}^{L} \sin \left( \frac{n\pi x}{\ell } \right) \sin \left( \frac{m\pi x}{\ell } \right)  \,\mathrm{d} x\right)=\int_{0}^{L} u_0 \sin \left( \frac{m\pi x}{\ell } \right)  \,\mathrm{d} x
.\]
The left integral is the inner product of our eigenfunctions, so it follows by orthogonality that the integral is $0$ for all $n\neq m$, so we have
\[
	B_m \int_{0}^{\ell }\left(\sin \left( \frac{m\pi x}{\ell } \right)\right)^2 \mathrm{d} x=\int_{0}^{\ell } u_0\sin \left( \frac{m\pi x}{\ell } \right)  \,\mathrm{d} x
.\]
Therefore,
\[
	B_m =\frac{\int_{0}^{\ell } u_0\sin \left( \frac{m\pi x}{\ell } \right)  \,\mathrm{d} x}{\int_{0}^{\ell } \left( \sin \left( \frac{m\pi x}{\ell } \right)  \right) ^2 \,\mathrm{d} x}
.\]
We can simplify this formula somewhat. Using Euler's formula, $\left( \sin \theta  \right) ^2=\frac{1}{2}-\frac{1}{2}\cos (2\theta )$. It follows that
\[
	\int_{0}^{\ell } \frac{1}{2}-\frac{1}{2}\cos \left( \frac{2n\pi x}{\ell } \right)  \,\mathrm{d} x
.\]
Notice that each cosine is $0 $ after the integral, so we get $\ell  /2$. Therefore,
\[
	B_m=\frac{2}{\ell }\int_{0}^{\ell } u_0 \sin \left( \frac{n\pi x}{\ell } \right)  \,\mathrm{d} x
.\]
Let $u(x,t)$ be some function. Consider the partial differential equation
\[
	\frac{\partial u}{\partial t} =\frac{\partial ^2u}{\partial x^2} ,
\]
with initial condition $u(x,0)=u_0$ and boundary conditions $u(0,t)=u(\ell ,t)=0$. The general solution to this system is
\[
	u(x,t)=\sum_{n=1}^{\infty} \exp \left( -\left( \frac{2n\pi }{\ell } \right) ^2t \right) \frac{4u_0}{\ell }\sin \left( \frac{2n\pi }{\ell }x \right) 
.\]
\[
	u(x,t)=\sum_{n=1}^{\infty} \frac{4u_0 \sin \left( \frac{2n\pi }{\ell }x \right) }{\exp \left( \left( \frac{2n\pi }{\ell } \right)^2 t  \right) \ell }
.\]


